\subsection{Endogenous system characteristics}\label{\refsec:ESC}
This section documents how the design choices of \cref{\refmodelsec:ESC} are solved. Two solvers are described, one per timing --- the sequential choice needs none, since \eqref{\refmodeleq:esc:seqFOC} is a scalar first order condition of the same form as the tax one --- together with the calibration of the cost parameter and the checks that establish the structural claims rather than assuming them.

Throughout, $\Theta = \left\lbrace \theta^{(1)},\dots,\theta^{(M_{\theta})}\right\rbrace$ denotes a grid on $[0,1]$ and $\mathcal{T}_t(\theta)$ the ordinary equilibrium tax at date $t$ given a design $\theta$, i.e.\ the solution of $z_t(\tau,\theta)=0$ selected as in \eqref{\refeq:candidates}. By the decoupling property this is a \emph{static} scalar problem, one root per node, with no dependence on $s_{t-1}$ and no recursion --- which is what makes the log case cheap.

\smalltitle{Why a grid and not a first order condition.} The question being asked is whether the choice is interior. A solver built on $\mathrm{d}\mathcal{W}/\mathrm{d}\theta = 0$ answers it only by failing to converge, and a corner is then indistinguishable from a badly behaved residual. We therefore evaluate the objective itself on $\Theta$, take the maximiser, and refine it by fitting a parabola through the three nodes around it when the maximum is interior. A corner is returned \emph{as} a corner. The cost is one objective evaluation per node; the benefit is that the central result of the exercise is observed rather than inferred. This inverts the convention used for the tax: there the first order condition is the cheap object and the objective is reconstructed from it (\cref{\refsec:RobustRoot}); here the objective is evaluated directly, precisely because its corner behaviour is the substance of the question.

\smalltitle{Evaluating the objective does not relax the discipline on predetermined states.} Direct evaluation may look like it surrenders what the first order condition was built to guarantee --- that the electorate takes the inherited distribution of savings as given. It does not, because that guarantee never resided in the representation. The tax residual $z_t$ excludes $s_{t-1,i}/s_{t-1}$ from the variation because it is \emph{constructed} as the derivative holding the ratio fixed; the grid implements the same definition by evaluating every candidate at the same predetermined ratio, and the two representations enforce one discipline. Under the leaded timing the discipline is automatic: the candidate $\theta_{t+1}$ does not appear in \eqref{\refmodeleq:esc:auxiliary:si} at vintage $t-1$, so no channel exists through which the sweep could move the ratio. What \emph{does} move with the candidate --- the current young's labour, savings and their distribution --- is meant to: those choices are being made at $t$, not sunk. Under the permanent timing the candidate does enter the ratio's formula, through $\theta_{t_0}$, and there the ratio is pinned explicitly; \cref{\refsec:ESC:permanent} states the pinning and measures what ignoring it would cost.


\subsubsection{The leaded choice}\label{\refsec:ESC:leaded}
The design chosen at $t$ is a state at $t+1$, so the object to identify is a pair of sequences of policy functions, $\mathcal{T}_t(\theta_t)$ and $\Theta_t(\theta_t)$. \Cref{prop:esc:separability} says the second is constant in its argument, but the algorithm does not use that: it tabulates the policy on $\Theta$ precisely so that the claim can be measured.

\begin{alg}\label{\refalg:esc:leaded}
\emph{Leaded choice, log preferences.} Iterate backwards over $t = T-1,\dots,1$.
\begin{enumerate}
    \item \emph{Terminal period.} Solve $\mathcal{T}_{T-1}(\theta)$ for every $\theta\in\Theta$ from the terminal first order condition, i.e.\ \eqref{\refmodeleq:PEELOG} evaluated with $\beta_{T,j}=0$. There is no $\theta_T$ to choose.
    \item \emph{Taxes.} At $t<T-1$, evaluate $z_t$ on the flattened mesh $\Theta\times\tilde{\mathcal{T}}$ in a single call --- $z_t$ is elementwise in both arguments, so the mesh costs one evaluation rather than $M_{\theta}$ --- and select the maximiser at each node as in \eqref{\refeq:candidates}. \emph{Polish} each interior selection by a bracketed scalar root find within one grid cell.
    \item \emph{Design.} For each $\theta_t\in\Theta$, evaluate
    \begin{align}\label{\refeq:esc:leadedObjective}
        \mathcal{W}_t\left(\theta';\theta_t\right) = \sum_i \gamma_{t-1,i}\omega p_{t-1}\mu_{t-1,i}\ln\left(c_{2,t}^i\right) + \nu_t\sum_i \gamma_{t,i}\mu_{t,i}\left[(1+\beta_{t,i})\ln\left(\tilde{c}_{1,t}^i\right)+\beta_{t,i}\ln\left(R_{t+1}\right)\right]
    \end{align}
    at every candidate $\theta'\in\Theta$, using $\tau_t = \mathcal{T}_t(\theta_t)$ from step 2 and the continuation $\tau_{t+1} = \mathcal{T}_{t+1}(\theta')$, $\theta_{t+2} = \Theta_{t+1}(\theta')$, $\tau_{t+2} = \mathcal{T}_{t+2}(\theta_{t+2})$ from the periods already solved. Set $\Theta_t(\theta_t)$ to the refined maximiser.
\end{enumerate}
\end{alg}

Three implementation points are worth recording, because each guards against a plausible wrong answer.

\emph{The polish is not optional.} The grid selection alone is $O(\text{spacing}^2)$ accurate, roughly $10^{-4}$ at the resolutions used, while $\tau_{t_0}$ is a calibration target matched at $10^{-8}$. Without the refinement the cost parameter would be calibrated against grid noise. Corners are never polished: $z_t \neq 0$ there is the correct answer, not a residual.

\emph{$s_{t-1}$ is set to one.} By \cref{prop:esc:separability} the maximiser does not depend on it. The normalisation is convenient rather than necessary and is verified numerically, by re-maximising \eqref{\refeq:esc:leadedObjective} at savings levels spanning two orders of magnitude and confirming the argument does not move.

\emph{Interpolation is piecewise linear throughout.} A policy function that is about to be maximised over must not carry interpolation artefacts: a smoother or a cubic can manufacture a spurious interior maximum out of a profile that is genuinely flat at a bound, and an argmax --- unlike a root --- has no residual through which such an artefact would announce itself. The cautions of \cref{\refsec:PEE} about interpolation kind therefore apply with more force here than in their original setting.


\subsubsection{The permanent choice}\label{\refsec:ESC:permanent}
Here the design is chosen once, at $t_0$, and the electorate chooses $\tau_{t_0}$ alongside it. \Cref{prop:esc:concentration} reduces what looks like a two-dimensional grid search over $\left(\tau_{t_0},\theta\right)$ to a one-dimensional one, and the reduction is what makes the exercise cheap:

\begin{alg}\label{\refalg:esc:permanent}
\emph{Permanent choice.} Given a pinned ratio $\bar{\sigma}_i$, for each $\theta\in\Theta$: (i) obtain $\tau_t = \mathcal{T}_t(\theta)$ for every $t\geq t_0$ --- for $t>t_0$ this is the ordinary equilibrium at a constant design, with no recursion in $\theta$ to solve; (ii) evaluate \eqref{\refeq:esc:leadedObjective} at $t_0$ with $\theta_{t_0}=\theta'=\theta$ and $s_{t_0-1,i}/s_{t_0-1}=\bar{\sigma}_i$. Take the refined maximiser over $\Theta$. The vote being anticipated, $\bar{\sigma}_i$ is itself the ratio implied by the winning design, so the design is the fixed point of
\begin{align}\label{\refeq:esc:permFixedPoint}
    \theta^{\ast} = \arg\max_{\theta\in\Theta}\ \mathcal{W}_{t_0}\left(\theta;\ \sigma_i(\theta^{\ast})\right),
\end{align}
reached by iterating the above from the incumbent design. Convergence is reported, not assumed: this is a best response and not a contraction, so a cycle has to be visible.
\end{alg}

\smalltitle{The predetermined ratio must be held fixed, and at which value.} This is the one place where the permanent timing differs from the leaded one in substance rather than in bookkeeping, and it raises two questions that are easily run together.

The first is whether the ratio may move with the candidate at all. It may not. Under the leaded timing $\theta_{t+1}$ does not appear in $s_{t_0-1,i}/s_{t_0-1}$ and the question never arises; under the permanent timing $\theta$ does appear there, through $\theta_{t_0}$. But savings made at $t_0-1$ are sunk when the vote is taken at $t_0$: a voter comparing two candidates does not face a different $s_{t_0-1,i}$ under each. Maximising $\mathcal{W}_{t_0}$ while recomputing the ratio at every candidate would credit the electorate with internalising a state it takes as given --- the same error that \cref{\refsec:PEE} forbids for $\tau_t$, arriving in a new guise --- and it is not a small one: at $p=0.4$ the two readings give $0.775$ and $0.910$. It is also inconsistent with the tax it would be paired with, since $s_{t_0-1,i}/s_{t_0-1}$ is a function of $\tau_{t_0}$ as well as $\theta_{t_0}$ while $\mathcal{T}_{t_0}$ solves $z_{t_0}=0$, which is constructed holding the ratio fixed. It would take \cref{prop:esc:concentration} with it.

The second question is which value the pinned ratio takes, and there the answer is the timing's. The reform is \emph{anticipated}: households arrive at $t_0$ knowing a design will be chosen, and what they choose among is which of the $\left|\Theta\right|$ equilibrium paths the economy follows from $t_0$ on. The savings they made at $t_0-1$ were therefore made against the design that wins, which is \eqref{\refeq:esc:permFixedPoint}. The incumbent's ratio --- the natural pinning if the reform were a surprise --- is the seed of that iteration and not its answer. The two coincide \emph{exactly} wherever the chosen design reproduces the incumbent one, which is precisely the calibration condition \eqref{\refeq:esc:calibration}, so the calibrated cost is common to both readings and only the counterfactuals separate them; at $p=0.4$, away from the permanent timing's own calibration, the anticipated reading gives $0.775$ against the unanticipated $0.773$, and at $p=0.25$ it gives $0.542$ against $0.549$. Both are reported.

That the fixed point is a scalar problem at all is worth one line. The ratio in \eqref{\refmodeleq:esc:auxiliary:si} at vintage $t_0-1$ is a function of date-$t_0$ policy alone, so the kinked path the reform actually is --- the incumbent design before $t_0$, the chosen one from $t_0$ --- delivers the same $\bar{\sigma}_i$ as a design that had always been $\theta$, and no separate solve is needed to evaluate a candidate's savings. The kink is nonetheless carried in the reported equilibrium path, where a design applied to periods before the vote would be a counterfactual past.

\smalltitle{Verification.} \Cref{prop:esc:concentration} is checked rather than trusted: the tax at the chosen design must equal $\mathcal{T}_{t_0}$ evaluated there, and does, to $10^{-6}$. Were the concentration to fail --- because some channel made $\theta$ a function of $\tau$ --- \cref{\refalg:esc:permanent} would be solving the wrong problem, and this is the check that would say so.


\subsubsection{CRRA preferences}\label{\refsec:ESC:crra}
None of the log simplifications survive. The powers in $\left(\tilde{c}\right)^{1-1/\rho}$ do not turn products into sums, so \eqref{\refmodeleq:esc:separability} fails, both $s_{t-1}$ and $\theta_t$ re-enter the design choice, and the Markov object is a policy function over the two-dimensional state $\left(s_{t-1},\theta_t\right)$ --- a second dimension on top of the grid machinery of \cref{\refsec:PEE}.

For the permanent timing this costs nothing: a permanent design means a \emph{constant} $\theta$ path, each candidate is one ordinary CRRA equilibrium solve, and \cref{\refalg:esc:permanent} goes through unchanged. For the leaded timing we solve the equilibrium \emph{path} rather than the policy function:

\begin{alg}\label{\refalg:esc:crraPath}
\emph{Leaded choice, CRRA preferences.} Given a design path $\bm{\theta}$ with $\theta_t$ fixed at the incumbent value for $t\leq t_0$: sweep forward over $t \geq t_0$; at each $t$, evaluate the CRRA analogue of \eqref{\refeq:esc:leadedObjective} at every candidate $\theta_{t+1}\in\Theta$, re-solving the whole politico-economic equilibrium for each candidate while holding $\theta_{t+2},\dots$ at the current iterate; replace $\theta_{t+1}$ by the maximiser. Repeat the sweep until the path stops moving.
\end{alg}

Re-solving at every candidate is what makes the envelope logic right: $\tau_t$ is re-optimised at each design, so its own response to $\theta_{t+1}$ contributes nothing to first order and the objective is evaluated at the equilibrium tax rather than a stale one.

\smalltitle{What \cref{\refalg:esc:crraPath} assumes, and how far it is true.} Holding $\theta_{t+2}$ fixed while $\theta_{t+1}$ varies is exact if the choice at $t+1$ does not respond to the design it inherits. Under log preferences that is \cref{prop:esc:separability} and no approximation at all. Under CRRA it is an approximation, and we measure it directly: perturbing $\theta_{t+1}$ and re-optimising $\theta_{t+2}$ at each perturbation gives $\mathrm{d}\theta_{t+2}/\mathrm{d}\theta_{t+1}$ of $-0.0085$ at $\rho=2$. The residual state dependence is two orders of magnitude below the movements being studied, so the path iteration is adequate here; at a calibration where it is not, the two-dimensional grid becomes necessary and this shortcut has to go.

\smalltitle{Validation against the log limit.} The CRRA solver has no closed form to check against, so it is held to continuity: as $\rho\rightarrow1$ its answer must converge on the log solver's, and at a first-order rate, since nothing in the problem is kinked at $\rho=1$. Holding the cost parameter and the calibration fixed, the gap is $0.0903$, $0.0448$ and $0.0188$ at $\rho = 1.10$, $1.05$ and $1.02$, so the ratio $\text{gap}/(\rho-1)$ is $0.90$, $0.90$, $0.94$. The two solvers share the weights of \eqref{\refeq:esc:leadedObjective} and nothing else --- one maximises a closed form on a state grid by backward recursion, the other iterates on a path of full equilibrium solves --- so their agreement in the limit is informative.

\smalltitle{The exact solution: policy functions over the two-dimensional state.} \Cref{\refalg:esc:crraPath} approximates the Markov object; the object itself is also computed, by extending the backward recursion of \cref{\refsec:PEE} with the design as a second state:

\begin{alg}\label{\refalg:esc:crra2D}
\emph{Leaded choice, CRRA preferences, exact.} Let $\mathcal{S}\times\Theta$ be a grid over the state $\left(s_{t-1},\theta_t\right)$ and $\Theta'$ a grid of candidate designs. Solve the terminal period as in \cref{\refsec:PEE} once per $\theta_T\in\Theta$. Then, backwards over $t$:
\begin{enumerate}
    \item For every $\left(\tau_t, s_{t-1}, \theta_{t+1}\right) \in \mathcal{T}\times\mathcal{S}\times\Theta'$, resolve the equilibrium state $s_t$ from the state residual of \cref{\refsec:PEE}, with the continuation $\left(\tau_{t+1}, h_{t+1}\right)$ read off the period-$(t+1)$ policy interpolants at $\left(s_t, \theta_{t+1}\right)$.
    \item Assemble $z_t$ on that grid. The inherited design $\theta_t$ enters only the current old's $\mathrm{d}\upsilon_{2,t}^i/\mathrm{d}\tau_t$ term --- it splits benefits already granted --- so the numerical $\tau$-derivatives are shared across the $\theta_t$ axis and only that one term is recomputed per node.
    \item Select $\tau^{\ast}\left(s_{t-1},\theta_t,\theta_{t+1}\right)$ by the downward-crossing criterion of \cref{\refsec:RobustRoot}, evaluate the objective at the selection, and take $\Theta_t\left(s_{t-1},\theta_t\right)$ as its maximiser over $\Theta'$, refined parabolically when interior; corners are reported as corners.
    \item Tabulate $\left(\tau_t, \theta_{t+1}, s_t, h_t\right)$ at the chosen policies and rebuild the interpolants for period $t-1$.
\end{enumerate}
\end{alg}

One backward pass --- no iteration and no convergence tolerance; periods at which the design is history rather than a choice (every $t\leq t_0$ in the baseline timing) are solved with the candidate set collapsed to the inherited design, which matters here because $\tau_t$ genuinely responds to $\theta_{t+1}$ under CRRA. Pinned everywhere, the recursion collapses to the exogenous-design solver of \cref{\refsec:PEE} on a redundant grid, and reproduces its tax path to $\max_t\left|\Delta\tau_t\right| = 3.2\times10^{-5}$ --- a check of every ingredient except the choice layer against the production solver. At the calibrated cost at $\rho=2$, the exact choice at $t_0$ is $0.759$ against the path iteration's $0.761$, a gap of $0.002$, and the tax at $t_0$ still lands on its target (drifts of $3.2\times10^{-4}$ in $\tau$ and $3.1\times10^{-4}$ in $R$). The objective is flat enough near its maximum that either method identifies the design only to about $\pm 0.01$, and the two agree comfortably within that band. The discretisation requirements are properties of the grids rather than of the cost parameter: the state grid is immaterial (the choice moves by $4\times10^{-4}$ between $50$ and $150$ nodes) while the design-\emph{state} grid is not ($7$ nodes misplace the choice by $0.02$; $13$ suffice). \Cref{\refalg:esc:crraPath} is therefore confirmed adequate for the counterfactual tables, with \cref{\refalg:esc:crra2D} as the standard it is held to.


\subsubsection{Calibrating the cost}\label{\refsec:ESC:calibration}
Three targets and three parameters. $\left(\beta,\omega\right)$ are calibrated to $\left(R_{t_0},\tau_{t_0}\right)$ exactly as in \cref{\refsec:calibration}, with the cost installed and the design held at $\theta^{\ast}$; $\phi$ is imposed; and $p$ is chosen so that the design \emph{in force} at the baseline year, on a path where the political choice binds from the first period of the horizon, is the observed one,
\begin{align}\label{\refeq:esc:calibration}
    \theta_{t_0} = \Theta_{t_0-1}\left(\theta_{t_0-1}\right) = \theta^{\ast}.
\end{align}
Read in words: the model's own political history delivers the observed system by the baseline year. Everything after $t_0$ is then a prediction, and the drift of the design as $\nu_t$ falls is the object of interest.

\smalltitle{Why the design in force and not the choice made.} The design is a state chosen one period earlier, so $\theta_{t_0}$ and $\Theta_{t_0}(\theta^{\ast})$ are different numbers: the first was picked under $\nu_{t_0-1}$, the second binds only at $t_0+1$. They would coincide if the policy function were time invariant, and $\nu_t$ is not. The distinction is forced by what the counterfactuals report. Those are equilibrium paths read at $t_0$ (\cref{\refsec:ESC:counterfactuals}), so the baseline path must sit on the observed design \emph{at} $t_0$ or every comparison in the tables is drawn against a number the model missed. Under log preferences \cref{prop:esc:separability} makes $\Theta_{t_0-1}$ constant in its argument, so \eqref{\refeq:esc:calibration} does not depend on the inherited $\theta_{t_0-1}$ and no fixed point in the design's own history arises; under CRRA the residual is evaluated with the rest of the path held at $\theta^{\ast}$, exactly as elsewhere in this section, which is again exact at the root. The two targets are close --- calibrating on the choice at $t_0$ instead gives $p = 0.402$ against $0.408$ under the proportional specification at $\rho=1$, and a design at $t_0$ of $0.727$ against the observed $0.738$ --- so what the choice of target buys is not a different economics but a baseline row that reproduces its own datum.

\smalltitle{The nesting is an approximation, and is reported as one.} $\left(\beta,\omega\right)$ are calibrated along the path with the design held fixed, not along the endogenous-design path. At the calibrated $p$ the two agree at $t_0$ and $t_0+1$ by construction and drift only later, so the error is second order --- but it is an error, and the solved endogenous path is therefore checked back against the $\left(R,\tau\right)$ targets it was calibrated to. The residuals are $9\times10^{-15}$ and $2.6\times10^{-9}$, which is what licenses treating the two blocks as separable.

This caveat belongs to the leaded timing alone. Under the permanent timing the chosen design is constant from $t_0$ on and, at the calibrated $p$, equal to the incumbent one --- so the exogenous path $\left(\beta,\omega\right)$ are calibrated along \emph{is} the endogenous path, period for period, and the nesting is an identity rather than an approximation.

\smalltitle{The permanent residual is the fixed-point condition.} Equation \eqref{\refeq:esc:calibration} asks for the $p$ at which the electorate re-elects the system it has, and under the permanent timing that is exactly \eqref{\refeq:esc:permFixedPoint}: a cost at which the choice taken against the incumbent's savings ratio reproduces the incumbent design is a cost at which the incumbent design is the fixed point. The two therefore share a root, and the residual is evaluated in the cheaper single-pass form; only away from the root do they differ, so a scan must not mix them.

\smalltitle{Why \eqref{\refeq:esc:calibration} is scanned and not bracketed.} The residual is \emph{flat at a corner}: wherever the choice sits at $\theta=0$ or $\theta=1$ it does not move with $p$ at all. A root finder handed such a bracket reports spurious convergence at whichever endpoint it happened to evaluate. We therefore scan $p$ on a logarithmic grid, locate a genuine sign change, and only then bracket. A run that finds no interior crossing reports that fact rather than returning a number --- which is the correct outcome for France, whose observed design is itself at the boundary. It is also what happens if the scanned interval is simply placed wrong: the required cost spans an order of magnitude across $\rho$, so an interval chosen at $\rho=2$ lies entirely below the root at $\rho=0.5$, and the scan says so instead of converging on its own endpoint. The interval scanned therefore spans all three.

\smalltitle{Two consequences of the identification.} Under the flat-only specification, \eqref{\refmodeleq:esc:RR} makes $\theta^{\ast}$ a function of $p$, so the target in \eqref{\refeq:esc:calibration} moves with the parameter being calibrated. This is handled by re-deriving $\theta^{\ast}$ from $RR_0$ at every trial $p$ --- the datum is the replacement-rate ratio, not the design --- and the inversion is bracketed by a scan, since $(1-\theta)/\theta$ is decreasing while $f$ is increasing and their product need not be monotone. Second, the required cost turns out to be nearly the same under all three timings, which is itself a useful check that the two solvers are solving the same economic problem.

\smalltitle{The calibrated cost across $\rho$.} Under CRRA the residual in \eqref{\refeq:esc:calibration} is evaluated with the future design held at $\theta^{\ast}$ while the single choice at $t_0$ varies --- exact at the root, where the held path \emph{is} the equilibrium path, and only an approximation away from it, where it merely has to locate the sign change. The calibrated cost falls steeply in the intertemporal elasticity: under the proportional specification $p = 0.965,\ 0.408,\ 0.090$ at $\rho = 0.5,\ 1,\ 2$ (and $1.461,\ 0.709,\ 0.178$ under the flat-only one, where $\theta^{\ast}$ moves with it: $0.690,\ 0.716,\ 0.733$). A higher elasticity strengthens the young's forward-looking stake in a Bismarckian design --- the same ordering the stake decomposition found --- so less friction is needed to keep the choice off the Beveridgean corner. All six points are in the repository's \texttt{results/esc/escCalibration\{,CRRA\}.csv}, produced by \texttt{python/US/runESC.py} and \texttt{runESCcrra.py} and consumed by the paper pipeline's \texttt{US\_ESC\_Calibration} table.


\subsubsection{The counterfactuals}\label{\refsec:ESC:counterfactuals}
Every counterfactual in this appendix --- and every one in the main text's US tables --- is a \emph{separate equilibrium path}, not a shock applied to the calibrated economy at the baseline year. The changed characteristics hold over the whole horizon, the economy begins at its own steady state rather than at the United States', the political choice binds from the first period, and the path is read at $t_0$. What a row describes is a country that has always had this mix of characteristics --- mostly American, partly French --- and the question the tables put is how far the observable characteristics carry the United States towards France.

\smalltitle{Why not an unanticipated reform at $t_0$.} An alternative construction dates the change at $t_0$, seeds the economy with the savings the baseline generation actually made, and so describes an economy surprised in 2020. That object is coherent, but it is not commensurable with France's own calibrated path, which is an equilibrium and not a surprise. Since the point of the French experiments is a comparison across countries, the paths have to be built the same way on both sides. The new-path construction has a second advantage: with the design pinned as history through $t_0$, a change dated $t_0$ could not move $\theta$ before $t_0+1$, so a table read at $t_0$ would be identical across the exogenous- and endogenous-design readings by construction. Here $\theta_{t_0}$ is itself an outcome, and 2020 carries the design response.

\smalltitle{What the convention costs, measured.} At $\rho=1$ the choice between the two constructions moves the average workweek and nothing else: the tax rate and the savings rate come out identical to every printed digit under both readings. Both are rate objects, and under log preferences with a Cobb--Douglas technology neither depends on the inherited capital stock, while the \emph{level} of hours responds to the wage and hence to $k_{t_0}$. This is a property of the log case and does not survive it --- under CRRA the tax responds to the state, and the $\rho\neq1$ rows move.

\smalltitle{Two readings.} Each counterfactual is reported both with the design held at the value the replacement-rate data imply under the changed characteristics, and with the design chosen. The two differ already at $t_0$, which is what makes a table read there informative. Alongside them the French tables carry France's own calibrated path. That row is not a counterfactual on the American model: France brings its own political weight $\omega$ as well as its own characteristics, so the distance between it and the all-characteristics row measures what the observables leave unexplained. Its workweek is a calibration target rather than a prediction, since the French calibration of \cref{\refsec:calibration} targets French hours relative to American ones, and is marked as such wherever it is printed.
